\section{Fermat varieties of degree \texorpdfstring{$114$}{114}}\label{sec:fermat114}

This section proves \Cref{thm:fermat114}. Aoki reduced the Hodge
conjecture for all Fermat varieties of a given degree $m$ to finitely many
classes, one for each divisor of $m$ of a certain kind
\cite{Aok83,Aok00}. For $m=114$ a single class remains, and we show that it
is algebraic on the Fermat eightfold (\Cref{ssec:aokigroups}). The class is
the join of two characters of the Fermat threefold
(\Cref{ssec:twochars}), and the main step is a proof that the two
corresponding pieces of $H^{3}$ are supported on divisors. For this we use
families of curves on the Fermat threefold and a residue formula for the
derivative of their Abel--Jacobi map. The method is Peterson's
\cite[\S3]{Pet26}, who used one such family on the Fermat threefold of
degree $35$; we set it up for families of a general shape
(\Cref{ssec:families}) and then write down two families of degree $114$
(\Cref{ssec:twofamilies}).

\subsection{Aoki's groups}\label{ssec:aokigroups}
In this subsection $M\ge3$ is an integer, and we use the language of
\Cref{ssec:fermatodd}: balanced multisets of nonzero elements of $\ZZ/M$,
pairs, and algebraic multisets. Let $R_{M}$ be the free abelian group on the
symbols $(y)$, $y\in\ZZ/M\setminus\{0\}$. A multiset $A$ of nonzero elements
has the class $[A]=\sum_{a\in A}(a)$, so that $[A+B]=[A]+[B]$, and an element
of $R_{M}$ with nonnegative coefficients is the class of exactly one
multiset. A unit $t$ acts by $t\cdot(y)=(ty)$. For $M=gM'$ the homomorphism
$g_{*}\colon R_{M'}\to R_{M}$, $(y)\mapsto(gy)$, maps $[A']$ to $[gA']$, in
the notation of \Cref{lem:algtools}(d).

Following Aoki \cite[\S2]{Aok00}, let $B_{M}\subseteq R_{M}$ be the subgroup
generated by the classes of the balanced multisets with an even number of
elements, let $D_{M}$ be the subgroup generated by the classes $(y)+(-y)$ of
the pairs, and let $S_{M}$ be the subgroup generated by $D_{M}$ and the
classes of the \emph{standard} multisets
\[
  \sigma_{p,i}=\Phi_{p}(i)+\{-pi\}\quad(p\ \text{odd}),\qquad
  \sigma_{2,i}=\{i,\ i+M/2,\ -2i,\ M/2\}.
\]
Here $p$ is a prime dividing $M$ with $p<M$, $d=M/p$, $0<i<d$, and
$\Phi_{p}(i)=\{i+jd:\ 0\le j<p\}$ as in \Cref{ssec:fermatodd}; for odd $p$,
$\sigma_{p,i}$ is the multiset $\sigma_{p,x}$ of that subsection with $x=i$.

\begin{lemma}\label{lem:standardalg}
Every standard multiset is a balanced multiset of an even number of nonzero
elements, and it is algebraic. Hence $D_{M}\subseteq S_{M}\subseteq B_{M}$.
These three groups are stable under the units, and for $M=gM'$ the map
$g_{*}$ carries each of them at level $M'$ into the same group at level $M$.
\end{lemma}

\begin{proof}
For $p=2$ the elements of $\sigma_{2,i}$ are nonzero because $0<i<M/2$, and
$\sigma_{2,i}$ is the multiset $T(i)$ of the proof of
\Cref{lem:exceptional}, which is balanced; having four elements, it is
algebraic by \Cref{lem:algtools}(a). Let $p$ be odd. Then $pi\ne0$, as
$0<i<d$, and the computation in the proof of \Cref{lem:algtools}(e), which
does not use that $M$ is odd, shows that $\sigma_{p,i}$ is a balanced
multiset of $p+1$ nonzero elements. For $p=3$ it is algebraic by
\Cref{lem:algtools}(a). For $p\ge5$ let $g=\gcd(i,d)$, $i'=i/g$, $d'=d/g$
and $M'=M/g$. Then $\sigma_{p,i}=g\,\sigma_{p,i'}$, where $\sigma_{p,i'}$ is
the standard multiset at level $M'$, $\gcd(i',d')=1$ and $d'\ge2$. If
$d'\ge3$, Aoki constructs a subvariety of $X^{p-1}_{M'}$ of codimension
$(p-1)/2$ that represents $\sigma_{p,i'}$ \cite[Theorem~2-1]{Aok87}. If
$d'=2$, then $i'=1$, $M'=2p$, and $\sigma_{p,1}=\{1,3,\dots,2p-1\}+\{p\}$ is
the union of the pairs $\{y,2p-y\}$ with $y<p$ odd and of the pair
$\{p,p\}$, so it is algebraic by \Cref{lem:algtools}(a),(b). In both cases
$\sigma_{p,i}$ is algebraic by \Cref{lem:algtools}(d).

For a unit $t$, multiplication by $t$ maps $\Phi_{p}(i)=i+d\ZZ/M$ onto
$\Phi_{p}(i'')$, where $0<i''<d$ and $i''\equiv ti\pmod d$, and it maps $-pi$
to $-pi''$, since $p(ti-i'')\in pd\ZZ=M\ZZ$; it fixes $M/2$ when $M$ is even,
as $t$ is then odd. So $t\sigma_{p,i}=\sigma_{p,i''}$. Similarly
$g\sigma_{p,i}$, formed at level $M'$, is $\sigma_{p,gi}$ at level $M$. As
balanced multisets and pairs are preserved by units and by $g_{*}$
(\Cref{lem:algtools}(c),(d)), the last assertions follow.
\end{proof}

\begin{lemma}\label{lem:transfer}
Let $\mathcal C$ be a set of algebraic balanced multisets, each with an even
number of elements, and let $A$ be a balanced multiset of an even number of
nonzero elements of $\ZZ/M$ with $[A]\in D_{M}+\sum_{C\in\mathcal C}\ZZ[C]$.
Then $A$ is algebraic. If $B_{M}=S_{M}+\sum_{C\in\mathcal C}\ZZ[C]$, then
every balanced multiset of an even number of nonzero elements of $\ZZ/M$ is
algebraic, and the Hodge conjecture holds for $X^{n}_{M}$ for every $n$.
\end{lemma}

\begin{proof}
Write $[A]=\sum_{l}\epsilon_{l}[C_{l}]+\sum_{y}k_{y}\bigl((y)+(-y)\bigr)$
with $C_{l}\in\mathcal C$, repetitions allowed, $\epsilon_{l}=\pm1$ and
$k_{y}\in\ZZ$. When $\epsilon_{l}=-1$ we use
$-[C_{l}]=[-C_{l}]-\sum_{c\in C_{l}}\bigl((c)+(-c)\bigr)$, where $-C_{l}$
is algebraic by \Cref{lem:algtools}(c). This gives
$[A]=\sum_{l}[C'_{l}]+\sum_{y}k'_{y}\bigl((y)+(-y)\bigr)$ with algebraic
$C'_{l}$ and integers $k'_{y}$. Let $P_{-}$ and $P_{+}$ be the unions of
$|k'_{y}|$ copies of the pair $\{y,-y\}$ over the $y$ with $k'_{y}<0$ and
with $k'_{y}>0$. Then $[A+P_{-}]=[C'_{1}+\dots+C'_{L}+P_{+}]$, so
$A+P_{-}=C'_{1}+\dots+C'_{L}+P_{+}$ as multisets. The right side is
algebraic by \Cref{lem:algtools}(a),(b), and so $A$ is algebraic by
\Cref{lem:algtools}(b).

If $B_{M}=S_{M}+\sum_{C}\ZZ[C]$, we apply this with the standard multisets
added to $\mathcal C$; they are algebraic by \Cref{lem:standardalg}. For
even $n$, the Hodge classes in $H^{n}(X^{n}_{M},\QQ)$ are spanned by a power
of the hyperplane class and the primitive Hodge classes, and by
\Cref{prop:fermatchars}(ii) the latter span, over $\CC$, the lines $V(a)$
whose entries form a balanced multiset; these lines are algebraic. In every
other degree $2k$ the group $H^{2k}(X^{n}_{M},\QQ)$ is spanned by a power of
the hyperplane class, by the Lefschetz hyperplane theorem and Poincar\'e
duality.
\end{proof}

\begin{lemma}\label{lem:nu}
Let $\ell$ be an odd prime with $\ell\parallel M$, and let $\nu_{\ell}\colon
R_{M}\to\ZZ/2$ be the homomorphism with $\nu_{\ell}((y))=1$ if $y$ has order
$\ell$ and $\nu_{\ell}((y))=0$ otherwise. Then $\nu_{\ell}$ vanishes on
$S_{M}$.
\end{lemma}

\begin{proof}
For a divisor $k$ of $M$ let $\Gamma_{k}\subseteq\ZZ/M$ be the subgroup of
order $k$. A pair has $0$ or $2$ elements of order $\ell$, since $y$ and
$-y$ have the same order. Let $\sigma=\sigma_{p,i}$ be standard and
$d=M/p$. The elements of $\sigma$ are those of the coset $i+\Gamma_{p}$,
which does not contain $0$, the element $-pi$, and, if $p=2$, the element
$M/2$, of order $2\ne\ell$.

Let $p\ne\ell$. Then $\Gamma_{p}\cap\Gamma_{\ell}=0$, so the coset
$i+\Gamma_{p}$ contains at most one element of $\Gamma_{\ell}$. It contains
one exactly when $pi\in\Gamma_{\ell}$: if $z\in(i+\Gamma_{p})\cap
\Gamma_{\ell}$ then $pi=pz$, and conversely, multiplication by $p$ is an
automorphism of $\Gamma_{\ell}$, so $pi=pz$ for some $z\in\Gamma_{\ell}$,
and then $i-z\in\Gamma_{p}$. In that case $-pi=-pz$ has the same order as
$z$. So $\sigma$ has $0$ or $2$ elements of order $\ell$.

Let $p=\ell$. Then $i\notin\Gamma_{\ell}=d\,\ZZ/M$, as $0<i<d$, so
$i+\Gamma_{\ell}$ does not meet $\Gamma_{\ell}$; and $-\ell i\in
\Gamma_{\ell}$ would mean that $d$ divides $\ell i$, hence $i$, since
$\gcd(\ell,d)=1$. So $\sigma$ has no element of order $\ell$.
\end{proof}

Let $\mathsf P$ be the set consisting of $4$ and the odd primes, and let
$\mathcal Q(M)$ be the set of divisors $q>1$ of $M$ that are products of an
even number of distinct elements of $\mathsf P$.

\begin{theorem}[Aoki]\label{thm:aoki}
For every $q\in\mathcal Q(M)$ there is an element $\xi_{q}\in B_{q}$ such that
\[
  B_{M}=S_{M}+\sum_{q\in\mathcal Q(M)}\ \sum_{t\in(\ZZ/M)^{\times}}
  \ZZ\,t\cdot(M/q)_{*}\xi_{q}.
\]
Moreover $2x\in S_{M}$ and $x+t\cdot x\in S_{M}$ for every $x\in B_{M}$ and
every unit $t$.
\end{theorem}

This is \cite[Theorem~2.1]{Aok00}, with the second assertion taken from the
proof of \cite[Corollary~2.3]{Aok00}; it rests on \cite[Theorem~D]{Aok83}.
In \cite{Aok83} the groups are normalized differently, and Peterson explains
the translation \cite[\S5.2]{Pet26}. Aoki also shows that
$B_{M}/S_{M}\cong(\ZZ/2)^{|\mathcal Q(M)|}$ \cite[Remark~2.4]{Aok00}; we
shall not need this.

\begin{corollary}\label{cor:B114}
Let $\beta$ be a balanced multiset of an even number of nonzero elements of
$\ZZ/114$ with $\nu_{19}([\beta])=1$. Then $B_{114}=S_{114}+\ZZ[\beta]$.
\end{corollary}

\begin{proof}
The elements of $\mathsf P$ dividing $114=2\cdot3\cdot19$ are $3$ and $19$,
so $\mathcal Q(114)=\{57\}$. By \Cref{thm:aoki} the group
$B_{114}/S_{114}$ is generated by the classes of the elements
$t\cdot2_{*}\xi_{57}$, and $t\cdot x\equiv-x\equiv x$ modulo $S_{114}$ for
every $x\in B_{114}$. So $B_{114}/S_{114}$ is generated by one element of
order at most $2$. By \Cref{lem:nu} with $\ell=19$, $[\beta]\notin S_{114}$;
hence $B_{114}/S_{114}$ has order $2$ and is generated by the class of
$[\beta]$.
\end{proof}

For the multiset $\beta$ of \Cref{lem:twochars} below, the equality
$B_{114}=S_{114}+\ZZ[\beta]$ has also been checked by an exact computation
with integer matrices that does not use \Cref{thm:aoki}
(\Cref{app:cert114}).

\subsection{Two characters of the Fermat threefold}\label{ssec:twochars}
Let $X=X^{3}_{m}$. Its cohomology $H^{3}(X,\QQ)$ is primitive, and by
\Cref{prop:fermatchars}(i), $H^{3}(X,\CC)$ is the direct sum of the lines
$V(a)$, $a\in\mathfrak A^{3}_{m}$, with $V(a)\subseteq H^{3-q,q}$ for
$q=\langle a\rangle-1$. As in the proof of \Cref{prop:fermatchars}(ii), for
$a\in\mathfrak A^{3}_{m}$ the sum $\bigoplus_{t}V(ta)$ over the units $t$ is
defined over $\QQ$; we write $M_{a}\subseteq H^{3}(X,\QQ)$ for the rational
sub-Hodge structure with $M_{a}\otimes\CC=\bigoplus_{t}V(ta)$. Since
$\langle-a\rangle=5-\langle a\rangle$, the structure $M_{a}$ has level one,
that is, $M_{a}\otimes\CC\subseteq H^{2,1}\oplus H^{1,2}$, exactly when
$\langle ta\rangle\in\{2,3\}$ for every unit $t$.

\begin{lemma}\label{lem:twochars}
Let $m=114$ and
\[
  \alpha_{1}=(1,25,43,57,102),\qquad \alpha_{2}=(39,63,68,80,92).
\]
\begin{enumerate}[label=\textup{(\alph*)}]
\item $\alpha_{1},\alpha_{2}\in\mathfrak A^{3}_{114}$, and for every unit $t$,
$\langle t\alpha_{1}\rangle,\langle t\alpha_{2}\rangle\in\{2,3\}$ and
$\langle t\alpha_{1}\rangle+\langle t\alpha_{2}\rangle=5$. In particular
$\langle\alpha_{1}\rangle=\langle13\alpha_{2}\rangle=2$, where
$13\alpha_{2}=(51,21,86,14,56)$.
\item The ten entries of $\alpha_{1}$ and $\alpha_{2}$ form a balanced
multiset $\beta$ with $\nu_{19}([\beta])=1$, and $\beta$ contains no pair.
\item $M_{\alpha_{1}}$ and $M_{\alpha_{2}}$ have dimension $36$.
\end{enumerate}
\end{lemma}

\begin{proof}
(a) The entries are nonzero, and they add up to $228$ and $342$. Since
$\langle-ta\rangle=5-\langle ta\rangle$, it suffices to consider the
eighteen units $t<57$, and for them the values are
\begin{center}
\small
\setlength{\tabcolsep}{3.2pt}
\begin{tabular}{l*{18}{c}}
\toprule
$t$ & 1 & 5 & 7 & 11 & 13 & 17 & 23 & 25 & 29 & 31 & 35 & 37 & 41 & 43 & 47 & 49 & 53 & 55\\
\midrule
$\langle t\alpha_{1}\rangle$ & 2&2&2&2&3&2&2&2&3&3&2&2&3&2&2&3&3&2\\
$\langle t\alpha_{2}\rangle$ & 3&3&3&3&2&3&3&3&2&2&3&3&2&3&3&2&2&3\\
\bottomrule
\end{tabular}
\end{center}
each entry being the sum of the five representatives in $\{1,\dots,113\}$
of the entries of $t\alpha_{i}$, divided by $114$. For instance, the entries
of $13\alpha_{2}=(51,21,86,14,56)$ add up to $228$.

(b) For every unit $t$ the representatives of the entries of $t\beta$ add up
to $114(\langle t\alpha_{1}\rangle+\langle t\alpha_{2}\rangle)=570$, so
$\beta$ is balanced. The elements of order $19$ in $\ZZ/114$ are the nonzero
multiples of $6$, and $102=6\cdot17$ is the only entry of this kind. No two
of the ten entries add up to $114$.

(c) If $t\alpha_{1}=\alpha_{1}$, then $t\equiv1$, since $1$ is an entry. If
$t\alpha_{2}=\alpha_{2}$, then $39t\equiv39$ and $68t\equiv68$, so $t\equiv1$
modulo $38$ and modulo $57$, hence modulo $114$. So the $36$ characters
$t\alpha_{1}$, and the $36$ characters $t\alpha_{2}$, are distinct.
\end{proof}

By (a), $V(\alpha_{1})$ and $V(13\alpha_{2})$ lie in $H^{2,1}(X^{3}_{114})$,
and the types of $M_{\alpha_{1}}$ and $M_{\alpha_{2}}$ are complementary:
$V(t\alpha_{1})\otimes V(t\alpha_{2})$ has type $(3,3)$ for every $t$. By
\Cref{lem:nu} and \Cref{cor:B114}, $\beta$ is not a combination of pairs and standard
multisets, so it is not reached by linear subspaces and Aoki's cycles. The
rest of the section shows that it is algebraic.

\subsection{Families of curves and coniveau}\label{ssec:families}
In this subsection $m\ge3$, $X=X^{3}_{m}\subseteq\PP^{4}$ is the Fermat
threefold with equation $F=\sum_{k}x_{k}^{m}$, $F_{k}=\partial F/\partial
x_{k}=m\,x_{k}^{m-1}$, $G=\mu_{m}^{5}/\mu_{m}$ acts on $X$ as in
\Cref{ssec:fermathodge}, and $U_{k}=X\cap\{x_{k}\ne0\}$, an affine open set.
We fix $\alpha\in\mathfrak A^{3}_{m}$ with $\langle\alpha\rangle=2$, and
write $\bar\alpha_{k}\in\{1,\dots,m-1\}$ for its entries, so that
$\sum_{k}\bar\alpha_{k}=2m$; then $V(\alpha)\subseteq H^{2,1}(X)$. Let
\[
  \Pi=\{z_{0}+\dots+z_{4}=0\}\subseteq\PP^{4},\qquad \Phi\colon X\to\Pi,\quad
  \Phi(x)=(x_{0}^{m}:\dots:x_{4}^{m}).
\]
The morphism $\Phi$ is finite and surjective, it is the quotient by $G$, and
it is \'etale over $\Pi^{\circ}=\Pi\cap\{z_{0}\cdots z_{4}\ne0\}$.

\begin{definition}\label{def:blockfamily}
Let $T$ be a smooth connected complex algebraic curve, not necessarily
complete. A \emph{block family of character $\alpha$ over $T$} consists of
binary forms $B_{1},\dots,B_{r}$ in $(s,t)$ of degrees $\delta_{1},\dots,
\delta_{r}\ge1$ and functions $c_{0},\dots,c_{4}$, all with coefficients in
the ring $\cO(T)$ of regular functions, and of integers $e_{kl}\ge0$
($0\le k\le4$, $1\le l\le r$), such that:
\begin{enumerate}[label=\textup{(B\arabic*)}]
\item for every $\tau\in T$ the forms $B_{l}(\tau)$ are nonzero and no two
of them have a common zero on $\PP^{1}$, and $c_{k}(\tau)\ne0$ for every
$k$;
\item $D=\sum_{l}e_{kl}\delta_{l}$ does not depend on $k$; for every $l$
some $e_{kl}$ is positive and some $e_{kl}$ vanishes; and the weight
$w_{l}=\sum_{k}\bar\alpha_{k}e_{kl}$ of every block is divisible by $m$;
\item the forms $u_{k}=c_{k}\prod_{l}B_{l}^{e_{kl}}$ of degree $D$ satisfy
$u_{0}+u_{1}+u_{2}+u_{3}+u_{4}=0$.
\end{enumerate}
We put $A=\prod_{l}B_{l}^{w_{l}/m}$, a form of degree $2D$, and
$\Psi=\prod_{k}c_{k}^{\bar\alpha_{k}}\in\cO(T)$; then
\begin{equation}\label{eq:uApsi}
  u_{0}^{\bar\alpha_{0}}u_{1}^{\bar\alpha_{1}}u_{2}^{\bar\alpha_{2}}
  u_{3}^{\bar\alpha_{3}}u_{4}^{\bar\alpha_{4}}=\Psi\,A^{m}.
\end{equation}
A \emph{chart} is an ordered pair $(i,j)$ of distinct indices such that no
block occurs in both $u_{i}$ and $u_{j}$, that is, $e_{il}e_{jl}=0$ for
every $l$, and such that no block occurring in $u_{i}$ vanishes at
$(s:t)=(1:0)$.
\end{definition}

Let $\tau\in T$. A common zero of the $u_{k}(\tau)$ would be a zero of some
$B_{l}$, by (B1), and $B_{l}$ is absent from some $u_{k}$ by (B2); so
$f_{\tau}=(u_{0}(\tau):\dots:u_{4}(\tau))$ is a morphism $\PP^{1}\to\Pi$, by
(B3). Let
\[
  \mathcal C_{\tau}=\PP^{1}\times_{\Pi}X=\{((s:t),x)\in\PP^{1}\times X:\
  f_{\tau}(s:t)=\Phi(x)\},
\]
a curve finite over $\PP^{1}$. For an irreducible component of
$\mathcal C_{\tau}$ we write $\widehat C$ for its normalization, and
$\pi_{1}\colon\widehat C\to\PP^{1}$ and $\nu\colon\widehat C\to X$ for the two
maps; $\pi_{1}$ is finite. Let $\widehat C^{\circ}$ be the preimage of the
set of points $(s:1)$ with $u_{0}\cdots u_{4}(s,1)\ne0$. Since $\Phi$ is
\'etale over $\Pi^{\circ}$, near every point of $\widehat C^{\circ}$ the
function $s$ is a local coordinate, and $\nu$ is $s\mapsto x(s,\tau)$ with
\begin{equation}\label{eq:liftgen}
  x(s,\tau')=\bigl(u_{0}(s,1;\tau')^{1/m},\dots,u_{4}(s,1;\tau')^{1/m}\bigr)
\end{equation}
for suitable branches, defined also for $\tau'$ near $\tau$.

\begin{definition}\label{def:residuefunctional}
Let $(i,j)$ be a chart, let $a<b<c$ be the other three indices, let
$\tau\in T$ and $v\in T_{\tau}T$. In the affine coordinate $s$ (that is,
$t=1$), denote by a dot the derivative along $v$ at fixed $s$ and by a prime
the derivative in $s$, and put
\[
  \Delta_{v}=\det\begin{pmatrix}1&1&1\\ \dot u_{a}/u_{a}&\dot u_{b}/u_{b}&
  \dot u_{c}/u_{c}\\ u'_{a}/u_{a}&u'_{b}/u_{b}&u'_{c}/u_{c}\end{pmatrix},
  \qquad \eta_{v}=\frac{A\,\Delta_{v}}{u_{i}u_{j}}\,ds,
\]
and
\[
  I_{ij}(\tau,v)=\sum_{\rho}\operatorname{Res}_{s=\rho}\eta_{v},
\]
the sum over the zeros $\rho\in\CC$ of $u_{i}(s,1)$. All zeros of $u_{i}$
lie in this affine line, because $(i,j)$ is a chart.
\end{definition}

\begin{lemma}\label{lem:etapoles}
\begin{enumerate}[label=\textup{(\alph*)}]
\item The poles of $\eta_{v}$ in $\CC$ lie at zeros of the blocks. At a zero
$\rho$ of $B_{l}$ of multiplicity $\mu$,
\[
  \operatorname{ord}_{\rho}\eta_{v}\ \ge\ \mu\Bigl(\frac{w_{l}}{m}-1-e_{il}
  -e_{jl}\Bigr).
\]
In particular $\eta_{v}$ is regular at the zeros of the blocks that occur
in neither $u_{i}$ nor $u_{j}$.
\item The map $(\tau,v)\mapsto I_{ij}(\tau,v)$ is a holomorphic $1$-form on
$T$.
\end{enumerate}
\end{lemma}

\begin{proof}
(a) We have $\dot u_{k}/u_{k}=\dot c_{k}/c_{k}+\sum_{l}e_{kl}\dot B_{l}/B_{l}$
and $u'_{k}/u_{k}=\sum_{l}e_{kl}B'_{l}/B_{l}$. Write $\mathbf1=(1,1,1)$,
$\gamma=(\dot c_{k}/c_{k})_{k=a,b,c}$ and $\epsilon_{l}=(e_{kl})_{k=a,b,c}$.
Expanding the determinant along its second and third rows,
\[
  \Delta_{v}=\sum_{l}\frac{B'_{l}}{B_{l}}\det(\mathbf1,\gamma,\epsilon_{l})
  +\sum_{l\ne l'}\frac{\dot B_{l}}{B_{l}}\frac{B'_{l'}}{B_{l'}}
  \det(\mathbf1,\epsilon_{l},\epsilon_{l'}),
\]
the terms with $l=l'$ vanishing. The functions $\dot c_{k}/c_{k}$ are
constants on $\PP^{1}$, by (B1). At a zero $\rho$ of $B_{l}$ of
multiplicity $\mu$ no other block vanishes, $B'_{l}/B_{l}$ has a simple pole
and $\dot B_{l}/B_{l}$ a pole of order at most $\mu$; so
$\operatorname{ord}_{\rho}\Delta_{v}\ge-\mu$. Moreover
$\operatorname{ord}_{\rho}A=\mu w_{l}/m$ and $\operatorname{ord}_{\rho}
(u_{i}u_{j})=\mu(e_{il}+e_{jl})$. The last assertion holds because
$w_{l}\ge m$, by (B2).

(b) Fix $\tau_{0}\in T$. Let $Z_{i}(\tau)\subseteq\CC$ be the set of zeros
of the blocks that occur in $u_{i}$, and $Z'(\tau)\subseteq\PP^{1}$ the set
of zeros of the other blocks. These sets are disjoint by (B1), and they
depend continuously on $\tau$, since the blocks occurring in $u_{i}$ do not
vanish at $(1:0)$. Choose disjoint closed discs around the points of
$Z_{i}(\tau_{0})$ that do not meet $Z'(\tau_{0})$, and let $\gamma$ be the
sum of their boundary circles. For $\tau$ near $\tau_{0}$ the discs still
contain $Z_{i}(\tau)$ and do not meet $Z'(\tau)$; so by (a)
$I_{ij}(\tau,v)=(2\pi i)^{-1}\oint_{\gamma}\eta_{v}$, which is holomorphic in
$(\tau,v)$ and linear in $v$.
\end{proof}

Peterson's argument now goes through for every block family; we give it in
full, with his notation where possible. Let $\xi=\sum_{k}x_{k}\partial_{k}$
be the Euler vector field on $\CC^{5}$, $\partial_{k}=\partial/\partial
x_{k}$, $dV=dx_{0}\wedge\dots\wedge dx_{4}$ and $\Omega=\iota_{\xi}dV$. For
$k\ne l$ put $\Omega_{kl}=\iota_{\partial_{l}}\iota_{\partial_{k}}\Omega$ and
\[
  \varphi_{kl}=\frac{x^{\bar\alpha-1}\,\Omega_{kl}}{F_{k}F_{l}},\qquad
  x^{\bar\alpha-1}=\prod_{n=0}^{4}x_{n}^{\bar\alpha_{n}-1}.
\]

\begin{lemma}\label{lem:cocyclegen}
\begin{enumerate}[label=\textup{(\alph*)}]
\item $\varphi_{kl}$ is the pullback of a rational $2$-form on $\PP^{4}$
that is regular where $x_{k}x_{l}\ne0$; its restriction to $X$ is a regular
$2$-form on $U_{k}\cap U_{l}$.
\item $\varphi_{ln}-\varphi_{kn}+\varphi_{kl}=0$ on $U_{k}\cap U_{l}\cap
U_{n}$, and $\varphi_{lk}=-\varphi_{kl}$.
\item $g^{*}\varphi_{kl}=\alpha(g)\varphi_{kl}$ for $g\in G$.
\end{enumerate}
So $\{\varphi_{kl}\}$ is a \v{C}ech $1$-cocycle of $\Omega^{2}_{X}$ for the
cover $\{U_{k}\}$, and its class $[\varphi]\in H^{1}(X,\Omega^{2}_{X})$ lies
in the $\alpha$-eigenspace.
\end{lemma}

\begin{proof}
(a) Under $x\mapsto\lambda x$ the factors $x^{\bar\alpha-1}$, $\Omega_{kl}$
and $F_{k}F_{l}$ are multiplied by $\lambda^{2m-5}$, $\lambda^{3}$ and
$\lambda^{2m-2}$, so $\varphi_{kl}$ is invariant; and $\iota_{\xi}
\Omega_{kl}=\iota_{\xi}\iota_{\partial_{l}}\iota_{\partial_{k}}\iota_{\xi}dV
=0$, since contractions anticommute. Hence $\varphi_{kl}$ descends to
$\PP^{4}$.

(b) The second assertion is clear. Since $\iota_{\xi}dF=mF$, we have
$dF\wedge\Omega=mF\,dV$. Contracting with $\partial_{k}$, $\partial_{l}$,
$\partial_{n}$, and using that $\iota_{\partial_{p}}$ is an antiderivation
with $\iota_{\partial_{p}}dF=F_{p}$, gives
\[
  F_{k}\Omega_{ln}-F_{l}\Omega_{kn}+F_{n}\Omega_{kl}
  =dF\wedge\iota_{\partial_{n}}\Omega_{kl}+mF\,\iota_{\partial_{n}}
  \iota_{\partial_{l}}\iota_{\partial_{k}}dV
\]
up to a sign common to both sides. On the cone over $X$ both $F$ and $dF$
restrict to zero; multiplying by $x^{\bar\alpha-1}/(F_{k}F_{l}F_{n})$ gives
the cocycle relation.

(c) For $g=(\zeta_{0}:\dots:\zeta_{4})$ we have $g^{*}x^{\bar\alpha-1}=
\prod_{n}\zeta_{n}^{\bar\alpha_{n}-1}x^{\bar\alpha-1}$ and $g^{*}F_{k}=
\zeta_{k}^{-1}F_{k}$, and $g^{*}\Omega_{kl}=\prod_{n\ne k,l}\zeta_{n}\,
\Omega_{kl}$, because every term of $\Omega_{kl}$ is a multiple of
$x_{n}\,dx_{n'}\wedge dx_{n''}$ with $\{n,n',n''\}=\{0,\dots,4\}\setminus
\{k,l\}$. The product of the factors is $\prod_{n}\zeta_{n}^{\bar\alpha_{n}}
=\alpha(g)$.
\end{proof}

This is the cocycle that Carlson and Griffiths attach to the residue of
$x^{\bar\alpha-1}\Omega/F^{2}$ \cite[\S3(b)]{CG80}; we do not need their
normalization. By \Cref{prop:fermatchars}(i), the $\alpha$-eigenspace of
$H^{2,1}(X)\cong H^{1}(X,\Omega^{2}_{X})$ is the line $V(\alpha)$.

We recall the \v{C}ech description of Dolbeault cohomology
\cite[Chapter~4]{Voi02b}. If $\theta$ is a $\bar\partial$-closed form of type
$(2,1)$ on $X$, there are smooth $(2,0)$-forms $\chi_{k}$ on the Stein open
sets $U_{k}$ with $\bar\partial\chi_{k}=\theta|_{U_{k}}$, the forms
$h_{kl}=\chi_{l}-\chi_{k}$ are holomorphic and form a \v{C}ech $1$-cocycle of
$\Omega^{2}_{X}$, and $\theta\mapsto\{h_{kl}\}$ induces the isomorphism
$H^{2,1}_{\bar\partial}(X)\cong H^{1}(X,\Omega^{2}_{X})$, compatibly with
the action of $G$.

\begin{definition}\label{def:firstorder}
Let $\tau\in T$, $v\in T_{\tau}T$, and let $\widehat C,\pi_{1},\nu$ be as
above. The vector $\partial x(s,\tau)/\partial v\in\CC^{5}$, the derivative
of \eqref{eq:liftgen} along $v$ at fixed $s$, is tangent to the cone over
$X$, since $\sum_{k}u_{k}=0$ identically; so it defines a tangent vector to
$X$ at $\nu(s)$. A \emph{first-order deformation of $\nu$ along $v$} is a
family of holomorphic sections $w_{q}$ of $\nu^{*}T_{X}$ over the members
$O_{q}$ of an open cover of $\widehat C$, such that $w_{q}-w_{q'}$ is
tangent to $\widehat C$ on $O_{q}\cap O_{q'}$ and $w_{q}\equiv\partial
x/\partial v$ modulo tangent vectors of $\widehat C$ on $O_{q}\cap
\widehat C^{\circ}$.
\end{definition}

For such $w$ and a smooth form $h$ of type $(2,0)$ or $(2,1)$ near
$\nu(\widehat C)$, the form $\nu^{*}(\iota_{w}h)$, defined locally as
$\nu^{*}(\iota_{w_{q}}h)$, does not depend on the choices: for $\varsigma$
tangent to $\widehat C$, $\nu^{*}(\iota_{d\nu(\varsigma)}h)=\iota_{\varsigma}
\nu^{*}h=0$, as $\nu^{*}h$ is a form of type $(2,0)$ or of degree three on a
curve. Since the $w_{q}$ are holomorphic, $\bar\partial\,\nu^{*}(\iota_{w}h)
=-\nu^{*}(\iota_{w}\bar\partial h)$ for $h$ of type $(2,0)$.

\begin{lemma}\label{lem:tracegen}
Let $(i,j)$ be a chart, $w$ a first-order deformation of $\nu$ along $v$,
and $W_{k}=\nu^{-1}(U_{k})$. For a \v{C}ech $1$-cocycle $h$ of
$\Omega^{2}_{X}$ on $\{U_{k}\}$ put
\[
  \mathcal R_{w}(h)=\sum_{o\in\widehat C\setminus W_{i}}\operatorname{Res}_{o}
  \nu^{*}(\iota_{w}h_{ij}).
\]
Then $\mathcal R_{w}(h)$ depends only on the class of $h$, and for every
$\bar\partial$-closed $(2,1)$-form $\theta$ with associated cocycle $h$,
\[
  \int_{\widehat C}\nu^{*}(\iota_{w}\theta)=-2\pi i\,\mathcal R_{w}(h).
\]
\end{lemma}

\begin{proof}
As $u_{i}$ and $u_{j}$ have no common zero, $x_{i}$ and $x_{j}$ do not
vanish together on $\nu(\widehat C)$; so $W_{i}\cup W_{j}=\widehat C$, and
$\widehat C\setminus W_{i}$ is finite. Put $g_{k}=\nu^{*}(\iota_{w}\chi_{k})$
on $W_{k}$ for $k=i,j$. Then $\bar\partial g_{k}=-\nu^{*}(\iota_{w}\theta)$,
and $g_{j}-g_{i}=\nu^{*}(\iota_{w}h_{ij})$ is holomorphic on $W_{i}\cap
W_{j}$. Choose small disjoint discs $D_{o}\subseteq W_{j}$ around the points
$o\in\widehat C\setminus W_{i}$. On a curve $\partial g_{k}=0$, so
$\nu^{*}(\iota_{w}\theta)=-dg_{k}$ on $W_{k}$, and Stokes' theorem gives
\[
  \int_{\widehat C}\nu^{*}(\iota_{w}\theta)
  =\sum_{o}\oint_{\partial D_{o}}g_{i}-\sum_{o}\oint_{\partial D_{o}}g_{j}
  =-2\pi i\sum_{o}\operatorname{Res}_{o}(g_{j}-g_{i}).
\]
If $h$ is a coboundary, $h_{kl}=\chi_{l}-\chi_{k}$ with holomorphic
$\chi_{k}$, and the same computation with $\theta=0$ gives
$\mathcal R_{w}(h)=0$.
\end{proof}

\begin{lemma}\label{lem:contractiongen}
Let $(i,j)$ be a chart. There is a sign $\epsilon_{ij}$, depending only on
$(i,j)$, such that on $\widehat C^{\circ}$, in the coordinate $s$,
\[
  \nu^{*}(\iota_{w}\varphi_{ij})=\epsilon_{ij}\,m^{-4}\,\nu^{*}\!\Bigl(
  \frac{x^{\bar\alpha}}{x_{i}^{m}x_{j}^{m}}\Bigr)\Delta_{v}\,ds,\qquad
  x^{\bar\alpha}=\prod_{n}x_{n}^{\bar\alpha_{n}}.
\]
Moreover there is $r\in\CC$ with $r^{m}=\Psi(\tau)$ such that
$\nu^{*}\bigl(x^{\bar\alpha}/(x_{i}^{m}x_{j}^{m})\bigr)=r\,\pi_{1}^{*}
\bigl(A/(u_{i}u_{j})\bigr)$ on $\widehat C$.
\end{lemma}

\begin{proof}
Work on the cone, with the lift \eqref{eq:liftgen}. The tangent vector and
the deformation vector are
\[
  \frac{\partial x}{\partial s}=\Bigl(\frac{x_{k}u'_{k}}{m\,u_{k}}\Bigr)_{k},
  \qquad\frac{\partial x}{\partial v}=\Bigl(\frac{x_{k}\dot u_{k}}{m\,u_{k}}
  \Bigr)_{k}.
\]
Since $\varphi_{ij}$ is the pullback of a form on $\PP^{4}$, its value on the
images of these vectors may be computed on the cone, and
$\Omega_{ij}(\partial x/\partial v,\partial x/\partial s)=dV(\xi,\partial_{i},
\partial_{j},\partial x/\partial v,\partial x/\partial s)$. Expanding this
$5\times5$ determinant along the rows of $\partial_{i}$ and $\partial_{j}$
leaves, up to a sign $\epsilon_{ij}$, the determinant of the rows $\xi$,
$\partial x/\partial v$, $\partial x/\partial s$ in the columns $a,b,c$,
which is $x_{a}x_{b}x_{c}m^{-2}\Delta_{v}$. Dividing by $F_{i}F_{j}=m^{2}
x_{i}^{m-1}x_{j}^{m-1}$ and multiplying by $x^{\bar\alpha-1}$ gives the
formula, since $w\equiv\partial x/\partial v$ modulo tangent vectors of
$\widehat C$.

Both sides of the second identity are rational functions on $\widehat C$,
of degree $0$ in $x$ and in $(s,t)$. On $\widehat C^{\circ}$, by
\eqref{eq:liftgen} and \eqref{eq:uApsi}, the $m$-th power of their quotient
is $\prod_{k}u_{k}^{\bar\alpha_{k}}/A^{m}=\Psi(\tau)\ne0$. So the quotient is
a constant on the connected curve $\widehat C$.
\end{proof}

\begin{proposition}[Residue formula]\label{prop:residuegen}
Let $\omega$ be the harmonic representative of a nonzero class in
$V(\alpha)$, and define $\lambda\in\CC$ by $[\varphi]=\lambda[\omega]$ in
$H^{1}(X,\Omega^{2}_{X})$. Let $(i,j)$ be a chart, $\tau\in T$, $v\in
T_{\tau}T$, $\widehat C$ the normalization of a component of
$\mathcal C_{\tau}$, and $w$ a first-order deformation of $\nu$ along $v$.
Then
\[
  \lambda\int_{\widehat C}\nu^{*}(\iota_{w}\omega)=-2\pi i\,\epsilon_{ij}\,
  m^{-4}\,r\,\deg(\pi_{1})\,I_{ij}(\tau,v),
\]
with $r$ as in \Cref{lem:contractiongen}. In particular, if $I_{ij}(\tau,v)
\ne0$, then $\lambda\ne0$ and $\int_{\widehat C}\nu^{*}(\iota_{w}\omega)\ne0$.
\end{proposition}

\begin{proof}
The class of $\omega$ spans the $\alpha$-eigenspace of $H^{1}(X,
\Omega^{2}_{X})$, which contains $[\varphi]$ by \Cref{lem:cocyclegen}; so
$\lambda$ is defined. Let $h$ be the cocycle attached to $\omega$. By
\Cref{lem:tracegen}, $\int_{\widehat C}\nu^{*}(\iota_{w}\omega)=-2\pi i\,
\mathcal R_{w}(h)$ and $\mathcal R_{w}(\varphi)=\lambda\mathcal R_{w}(h)$. By
\Cref{lem:contractiongen}, the holomorphic $1$-form $\nu^{*}(\iota_{w}
\varphi_{ij})$ on $W_{i}\cap W_{j}$ and the meromorphic $1$-form
$\epsilon_{ij}m^{-4}r\,\pi_{1}^{*}\eta_{v}$ agree on $\widehat C^{\circ}$, hence
on the connected set $W_{i}\cap W_{j}$, which contains a punctured
neighbourhood of every point of $\widehat C\setminus W_{i}$. These points are
the preimages of the zeros of $u_{i}$. For a meromorphic $1$-form $\eta$
near $\rho\in\PP^{1}$ one has $\sum_{o\in\pi_{1}^{-1}(\rho)}\operatorname{Res}_{o}
\pi_{1}^{*}\eta=\deg(\pi_{1})\operatorname{Res}_{\rho}\eta$. Therefore
$\mathcal R_{w}(\varphi)=\epsilon_{ij}m^{-4}r\deg(\pi_{1})I_{ij}(\tau,v)$,
and the formula follows. Finally $r\ne0$ and $\deg\pi_{1}\ge1$.
\end{proof}

Up to sign, $\int_{\widehat C}\nu^{*}(\iota_{w}\omega)$ is the derivative of
the Abel--Jacobi map of the family of curves, paired with $\omega$
\cite{CG80}; we only use it through the fibre integral below. Note that the
formula concerns one component at a time: the components of
$\mathcal C_{\tau}$ are permuted by $G$, and as $\alpha$ is a nontrivial
character, their contributions add up to zero.

\begin{lemma}\label{lem:familygen}
There are a smooth connected projective curve $B$ with a nonconstant
morphism $\pi\colon B\to\overline T$ to the smooth completion of $T$, a
smooth projective surface $S$ with morphisms $p\colon S\to B$ and $q\colon
S\to X$, and a nonempty Zariski-open set $B^{\circ}\subseteq B$, such that
for every $b\in B^{\circ}$:
\begin{enumerate}[label=\textup{(\roman*)}]
\item $\pi(b)\in T$, $\pi$ is \'etale at $b$, and $p$ is smooth over
$B^{\circ}$;
\item the fibre $S_{b}=p^{-1}(b)$, with the restriction of $q$ and of the
projection $S\to\PP^{1}$, is the normalization of an irreducible component
of $\mathcal C_{\pi(b)}$ with its maps $\nu$ and $\pi_{1}$;
\item for $v\in T_{b}B$, the images under $dq$ of local holomorphic lifts of
$v$ to $S$ form a first-order deformation of $\nu$ along $d\pi(v)$.
\end{enumerate}
\end{lemma}

\begin{proof}
Let $\mathcal C=\{(\tau,(s:t),x)\in T\times\PP^{1}\times X:\ f_{\tau}(s:t)=
\Phi(x)\}$. The morphism $\Phi$ is finite and flat, being a finite morphism
of smooth varieties of the same dimension, so $\mathcal C\to T\times\PP^{1}$
is finite and flat, and every irreducible component of $\mathcal C$ is a
surface mapping onto $T\times\PP^{1}$. Fix one, $\mathcal C'$. Let $B$ be
the smooth projective curve whose function field is the algebraic closure of
$\CC(T)$ in $\CC(\mathcal C')$, with the induced morphism $\pi$. The
inclusion of function fields gives a rational map $\mathcal C'
\dashrightarrow B$, and with the projections to $\PP^{1}$ and $X$ it maps
$\mathcal C'$ birationally onto a dense open subset of an irreducible
surface $\mathcal C''\subseteq B\times\PP^{1}\times X$, whose generic fibre
over $B$ is geometrically irreducible, since $\CC(B)$ is algebraically closed
in $\CC(\mathcal C'')$. Let $S$ be a resolution of singularities of
$\mathcal C''$ that is an isomorphism over its smooth locus, with the
induced $p$, $q$ and $S\to\PP^{1}$. By \cite[9.7.7]{EGA43} and generic
smoothness \cite[III.10.7]{Har77} there is a nonempty open $B^{\circ}$
satisfying (i) over which the fibres of $p$ are smooth, connected and
birational to the irreducible curves $\mathcal C''_{b}$. The image of
$\mathcal C''_{b}$ in $\PP^{1}\times X$ lies in $\mathcal C_{\pi(b)}$ and is
finite over $\PP^{1}$, so it is a component; this gives (ii). Over the open
set of $B\times\PP^{1}$ where $\pi$ is \'etale and $u_{0}\cdots u_{4}\ne0$,
the map $\mathcal C''\to B\times\PP^{1}$ is \'etale; so for a local
coordinate $z$ on $B$, the pair $(z,s)$ is a system of local coordinates on
$S$, and $q=x(s,\pi(z))$ with $x$ as in \eqref{eq:liftgen}. The vector field
with $s$-component $0$ and constant $z$-component equal to that of $v$ is
then a local lift $\tilde v$ of $v$ with $dq(\tilde v)=\partial x/\partial
(d\pi(v))$. Near the other points of $S_{b}$ lifts exist because $p$ is a
submersion, and two lifts differ by a vertical field, whose image under
$dq$ is tangent to the fibre. This gives (iii).
\end{proof}

\begin{theorem}\label{thm:coniveau}
Let a block family of character $\alpha$ over $T$ and a chart $(i,j)$ be
given, and suppose that $I_{ij}(\tau_{0},v_{0})\ne0$ for some $\tau_{0}\in
T$ and $v_{0}\in T_{\tau_{0}}T$. Then there are a smooth projective curve
$C$, possibly disconnected, and an algebraic cycle $Z$ of dimension $2$ on
$C\times X$, with the following property. Let $Z_{*}\colon H^{1}(C,\QQ)\to
H^{3}(X,\QQ)$ and $Z^{*}\colon H^{3}(X,\QQ)\to H^{1}(C,\QQ)$ be the induced
maps and $\Lambda=Z_{*}Z^{*}$. Then $\Lambda$ preserves $M_{\alpha}$, and
there is a polynomial $\ell$ with rational coefficients such that
$\ell(\Lambda)\Lambda$ is the identity on $M_{\alpha}$. In particular
$M_{\alpha}\subseteq Z_{*}H^{1}(C,\QQ)$, and $M_{\alpha}$ is supported on a
divisor of $X$.
\end{theorem}

\begin{proof}
Let $B,S,p,q$ be as in \Cref{lem:familygen} and $Z_{1}=(p,q)_{*}[S]$, a
cycle on $B\times X$, so that $Z_{1}^{*}=p_{*}q^{*}$. Let $\omega$ be the
harmonic representative of a nonzero class $\omega_{\alpha}\in V(\alpha)$, a
closed form of type $(2,1)$. We first show that $Z_{1}^{*}\omega_{\alpha}\ne
0$. Over $B^{\circ}$ the map $p$ is a proper submersion, and
$\vartheta=p_{*}q^{*}\omega$ is a closed $1$-form of type $(1,0)$ there,
hence holomorphic. For $b\in B^{\circ}$ and $v\in T_{b}B$, and any smooth
lift $\tilde v$ of $v$ along $S_{b}$, $\vartheta(v)=\pm\int_{S_{b}}
\iota_{\tilde v}q^{*}\omega$, and the integrand does not change if $\tilde v$
changes by a vertical field; so $\vartheta(v)=\pm\int_{S_{b}}\nu^{*}
(\iota_{w}\omega)$ with $w$ as in \Cref{lem:familygen}(iii). The form
$\vartheta$ extends holomorphically over each point $b_{0}\in B\setminus
B^{\circ}$: on a small punctured disc around $b_{0}$ its coefficient is
integrable, being bounded by the fibre integral of an integrable form, so it
has at most a simple pole; and the residue is $\pm(2\pi i)^{-1}
\int_{p^{-1}\{|z|\le\varepsilon\}}q^{*}d\omega=0$. As $p^{-1}(B\setminus
B^{\circ})$ has measure zero, $\vartheta$ represents $Z_{1}^{*}
\omega_{\alpha}$.

By \Cref{lem:etapoles}(b) the $1$-form $I_{ij}$ on the connected curve $T$ is
holomorphic and not identically zero, so it vanishes only on a discrete set,
while $\pi(B^{\circ})\cap T$ is $T$ minus a finite set. Choose $b\in
B^{\circ}$ with $I_{ij}(\pi(b),\cdot)\ne0$ and $v\in T_{b}B$ with $v\ne0$.
Since $\pi$ is \'etale at $b$, $I_{ij}(\pi(b),d\pi(v))\ne0$, and by
\Cref{prop:residuegen} and \Cref{lem:familygen}(ii),(iii), $\vartheta(v)\ne0$.
A nonzero holomorphic $1$-form on a compact Riemann surface is not exact, so
$Z_{1}^{*}\omega_{\alpha}=[\vartheta]\ne0$.

Now let $C=G\times B$, a disjoint union of $|G|$ copies of $B$, and let $Z$
restrict to $(\mathrm{id}_{B}\times g)_{*}Z_{1}$ on the copy indexed by
$g$. Then $\Lambda=\sum_{g\in G}g_{*}Z_{1*}Z_{1}^{*}g^{*}$ commutes with
the abelian group $G$, so it preserves each line $V(t\alpha)$ and acts on it
by a scalar $\Lambda_{t}$. Since $\Lambda$ is defined over $\QQ$, the
argument of \Cref{prop:fermatchars}(ii) shows that $\Lambda_{t}$ is the
image of $\Lambda_{1}\in\QQ(\mu_{m})$ under $\zeta\mapsto\zeta^{t}$. Using
$\int_{X}g_{*}x\cup x'=\int_{X}x\cup g^{*}x'$, $g^{*}\omega_{\alpha}=
\alpha(g)\omega_{\alpha}$ and $Z_{1}^{*}\bar\omega_{\alpha}=\bar\vartheta$,
\[
  \Lambda_{1}\int_{X}\omega_{\alpha}\cup\bar\omega_{\alpha}
  =\sum_{g\in G}|\alpha(g)|^{2}\int_{X}Z_{1*}Z_{1}^{*}\omega_{\alpha}\cup
  \bar\omega_{\alpha}=|G|\int_{B}\vartheta\wedge\bar\vartheta\ne0,
\]
as $i\int_{B}\vartheta\wedge\bar\vartheta>0$. So every $\Lambda_{t}$ is
nonzero, $\Lambda$ is invertible on $M_{\alpha}$, and by the
Cayley--Hamilton theorem its inverse there is $\ell(\Lambda)$ for a
polynomial $\ell$ with rational coefficients. Then $M_{\alpha}=
\Lambda(M_{\alpha})\subseteq Z_{*}H^{1}(C,\QQ)$, and $Z_{*}H^{1}(C,\QQ)$
vanishes on the complement of the image of the support of $Z$, a surface in
$X$.
\end{proof}

\subsection{The join of two threefold characters}\label{ssec:join}
We use the cohomological form of the inductive structure, due to Shioda
\cite[Theorem~II]{Shi79}; see also \cite{SK79}. Let $r,s\ge1$ and let $Z$ be
the blow-up of $X^{r}_{m}\times X^{s}_{m}$ along $Y\cong X^{r-1}_{m}\times
X^{s-1}_{m}$, with the morphism $\psi\colon Z\to X^{r+s}_{m}$ of
\Cref{prop:sk}. Let $E\subseteq Z$ be the exceptional divisor, $j\colon E\to
Z$ the inclusion and $\varpi\colon E\to Y$ the projection. Then
$f=\psi_{*}j_{*}\varpi^{*}\colon H^{r+s-2}(Y,\QQ)\to H^{r+s}(X^{r+s}_{m},\QQ)$
maps $V(a')\otimes V(b')$ onto $V(a'*b')$ for $a'\in\mathfrak A^{r-1}_{m}$
and $b'\in\mathfrak A^{s-1}_{m}$, where $a'*b'$ is the juxtaposed character.
Since $f$ is a composite of pullbacks and pushforwards along morphisms, it
maps algebraic classes to algebraic classes \cite[Chapter~19]{Ful98}.

\begin{proposition}\label{prop:join}
Let $\alpha,\alpha'\in\mathfrak A^{3}_{m}$ with $\langle t\alpha\rangle+
\langle t\alpha'\rangle=5$ for every unit $t$, and suppose that each of
$\alpha$ and $\alpha'$ has the property of \Cref{thm:coniveau}. Then the
multiset of the ten entries of $\alpha$ and $\alpha'$ is balanced and
algebraic.
\end{proposition}

\begin{proof}
Let $X=X^{3}_{m}$. The multiset is balanced because the entries of $t\alpha$
and $t\alpha'$ add up to $5m$. For every $t$, one of $V(t\alpha)$,
$V(t\alpha')$ lies in $H^{2,1}$ and the other in $H^{1,2}$; so $L=
\bigoplus_{t}V(t\alpha)\otimes V(t\alpha')\subseteq H^{6}(X\times X,\CC)$ has
type $(3,3)$. It is stable under $\operatorname{Gal}(\QQ(\mu_{m})/\QQ)$, so
$L=L_{\QQ}\otimes\CC$ for a space $L_{\QQ}\subseteq M_{\alpha}\otimes
M_{\alpha'}$ of Hodge classes.

Let $C,Z,\Lambda,\ell$ and $C',Z',\Lambda',\ell'$ be as in
\Cref{thm:coniveau} for $\alpha$ and $\alpha'$, and let $W=Z\times Z'$, a
correspondence from $C\times C'$ to $X\times X$. On $H^{3}(X)\otimes
H^{3}(X)$ the endomorphism $W_{*}W^{*}$ is $\pm\Lambda\otimes\Lambda'$, the
sign coming from the K\"unneth isomorphism, so $\pm\ell(\Lambda)\otimes
\ell'(\Lambda')$ inverts it on $M_{\alpha}\otimes M_{\alpha'}$. For
$\upsilon\in L_{\QQ}$ the class $\upsilon'=\pm W^{*}\bigl((\ell(\Lambda)
\otimes\ell'(\Lambda'))\upsilon\bigr)$ lies in $H^{2}(C\times C',\QQ)$ and is
a Hodge class, as $W^{*}$ and $\ell(\Lambda)\otimes\ell'(\Lambda')$ are
morphisms of Hodge structures up to a shift of type. By
\Cref{thm:lef11} on each component of the surface $C\times C'$,
$\upsilon'$ is algebraic, and so is $\upsilon=W_{*}\upsilon'$. Hence
$V(\alpha)\otimes V(\alpha')\subseteq L\subseteq\Alg^{3}(X\times X)\otimes
\CC$. By the inductive structure with $r=s=4$, the line $V(\alpha*\alpha')=
f\bigl(V(\alpha)\otimes V(\alpha')\bigr)$ is algebraic.
\end{proof}

\subsection{Two families}\label{ssec:twofamilies}
We now take $m=114$ and write down block families for the characters
$13\alpha_{2}$ and $\alpha_{1}$ of \Cref{lem:twochars}, both of which have
$\langle\cdot\rangle=2$. The first is explicit.

\begin{proposition}\label{prop:family1}
Let $\alpha=13\alpha_{2}=(51,21,86,14,56)$, $q(\lambda)=\lambda^{3}-
3\lambda^{2}+1$, and let $T_{1}\subseteq\CC$ be the complement of
$0,1,2/3,\pm1/\sqrt2$ and the roots of $q$. Over $T_{1}$, with parameter
$\lambda$, let the blocks be $s$, $s-t$, $s-\lambda t$ and $t$, and
\begin{align*}
  u_{0}&=q(\lambda)\,(s-\lambda t)\,t^{2}, & u_{1}&=-(\lambda-1)(s-\lambda
  t)^{3}, & u_{2}&=(\lambda-1)\,s\,(s-t)^{2},\\
  u_{3}&=-(\lambda-1)(3\lambda-2)\,s^{2}t, & u_{4}&=\lambda(2\lambda^{2}-1)
  (s-t)\,t^{2}.
\end{align*}
This is a block family of character $\alpha$ over $T_{1}$, $(2,0)$ is a
chart, and
\[
  I_{20}(\lambda,\partial_{\lambda})=-\frac{2\lambda^{3}+12\lambda^{2}
  -21\lambda+8}{\lambda(\lambda-1)^{2}(3\lambda-2)(2\lambda^{2}-1)\,
  q(\lambda)}.
\]
In particular $I_{20}(2,\partial_{\lambda})=5/28\ne0$.
\end{proposition}

\begin{proof}
On $T_{1}$ the blocks are pairwise coprime and the coefficients $c_{k}$ do
not vanish, which is (B1). The exponents of the blocks $s$, $s-t$,
$s-\lambda t$, $t$ in $(u_{0},\dots,u_{4})$ are $(0,0,1,2,0)$,
$(0,0,2,0,1)$, $(1,3,0,0,0)$ and $(2,0,0,1,2)$, with weights
$86+2\cdot14=114$, $2\cdot86+56=228$, $51+3\cdot21=114$ and
$2\cdot51+14+2\cdot56=228$; every $u_{k}$ is a cubic. This is (B2). For
(B3), in the coordinate $s$ ($t=1$) the coefficients of $s^{3}$, $s^{2}$,
$s$ and $1$ in $\sum_{k}u_{k}$ are
\begin{gather*}
  -(\lambda-1)+(\lambda-1),\qquad 3\lambda(\lambda-1)-2(\lambda-1)-(\lambda-1)
  (3\lambda-2),\\
  q(\lambda)-3\lambda^{2}(\lambda-1)+(\lambda-1)+\lambda(2\lambda^{2}-1),
  \qquad -\lambda q(\lambda)+\lambda^{3}(\lambda-1)-\lambda(2\lambda^{2}-1),
\end{gather*}
and all four vanish. The forms $u_{2}$ and $u_{0}$ share no block, and
$u_{2}(1,0)=\lambda-1\ne0$; so $(2,0)$ is a chart.

Here $A=s(s-t)^{2}(s-\lambda t)t^{2}$, so in the coordinate $s$,
$A/(u_{2}u_{0})=1/((\lambda-1)q(\lambda))$ is constant, and $\eta=\Delta\,ds/
((\lambda-1)q)$, where $\Delta$ is formed with the rows $1,3,4$ and $v=
\partial_{\lambda}$:
\begin{gather*}
  \frac{\dot u_{1}}{u_{1}}=\frac1{\lambda-1}-\frac3{s-\lambda},\qquad
  \frac{\dot u_{3}}{u_{3}}=\frac1{\lambda-1}+\frac3{3\lambda-2},\qquad
  \frac{\dot u_{4}}{u_{4}}=\frac1\lambda+\frac{4\lambda}{2\lambda^{2}-1}
  =\frac{6\lambda^{2}-1}{\lambda(2\lambda^{2}-1)},\\
  \frac{u'_{1}}{u_{1}}=\frac3{s-\lambda},\qquad\frac{u'_{3}}{u_{3}}=
  \frac2s,\qquad\frac{u'_{4}}{u_{4}}=\frac1{s-1}.
\end{gather*}
The zeros of $u_{2}$ are $s=0$ and $s=1$. At $s=0$ only $u'_{3}/u_{3}$ has a
pole, and expanding the determinant along its third row shows that the
residue of $\Delta$ is $2\bigl(\dot u_{1}/u_{1}(0)-\dot u_{4}/u_{4}\bigr)=
2\bigl(\frac1{\lambda-1}+\frac3\lambda-\frac{6\lambda^{2}-1}{\lambda
(2\lambda^{2}-1)}\bigr)$; over the common denominator $\lambda(\lambda-1)
(2\lambda^{2}-1)$ the numerator is $\lambda(2\lambda^{2}-1)+3(\lambda-1)
(2\lambda^{2}-1)-(\lambda-1)(6\lambda^{2}-1)=2\lambda^{3}-3\lambda+2$. At
$s=1$ only $u'_{4}/u_{4}$ has a pole, and the residue of $\Delta$ is
$\dot u_{3}/u_{3}-\dot u_{1}/u_{1}(1)=\frac1{\lambda-1}+\frac3{3\lambda-2}
-\frac4{\lambda-1}=-\frac{3(2\lambda-1)}{(\lambda-1)(3\lambda-2)}$. Hence
\[
  I_{20}(\lambda,\partial_{\lambda})=\frac{2(2\lambda^{3}-3\lambda+2)(3\lambda
  -2)-3(2\lambda-1)\lambda(2\lambda^{2}-1)}{\lambda(\lambda-1)^{2}(3\lambda-2)
  (2\lambda^{2}-1)\,q(\lambda)},
\]
and the numerator is $-(2\lambda^{3}+12\lambda^{2}-21\lambda+8)$. At
$\lambda=2$ the numerator is $-30$ and the denominator is
$2\cdot1\cdot4\cdot7\cdot(-3)=-168$.
\end{proof}

The second family has degree $24$, and we know of no family of smaller
degree for the orbit of $\alpha_{1}$. Let $(f,h,g,a,b,c,\kappa)$, with
$f=(f_{0},\dots,f_{6})$, $h=(h_{0},h_{1})$ and $g=(g_{0},\dots,g_{11})$, be
the coordinates of $\CC^{25}$, and consider the forms
\[
  F=s^{7}+\sum_{n<7}f_{n}s^{n}t^{7-n},\quad H=s^{2}+h_{1}st+h_{0}t^{2},\quad
  G=s^{12}+\sum_{n<12}g_{n}s^{n}t^{12-n},\quad K=s(s-t)(s-\kappa t),
\]
and
\[
  u_{0}=F^{3}K,\quad u_{1}=a\,HFK^{5},\quad u_{2}=b\,H^{5}F^{2},\quad
  u_{3}=-(1+a+b)\,G^{2},\quad u_{4}=c\,HK\,t^{19},
\]
binary forms of degree $24$. The coefficient of $s^{24}$ in
$E=\sum_{k}u_{k}$ is $1+a+b-(1+a+b)=0$; let $E_{0},\dots,E_{23}$ be the
other coefficients, $\mathcal P\subseteq\CC^{25}$ their common zero set, and
$J$ the Jacobian matrix of $(E_{0},\dots,E_{23})$ with respect to the
$24$ coordinates other than $\kappa$.

\begin{proposition}\label{prop:family2}
There is a point $p^{*}\in\mathcal P$ with $\kappa(p^{*})=(-17+3i)/10$ at
which
\begin{enumerate}[label=\textup{(\roman*)}]
\item $J$ is invertible;
\item $abc(1+a+b)\ne0$, and $F,H,G,K$ are pairwise coprime;
\item the residue functional of the chart $(0,3)$, defined below, does not
vanish on the tangent vector $v^{*}$ of $\mathcal P$ at $p^{*}$ with
$d\kappa(v^{*})=1$.
\end{enumerate}
\end{proposition}

The proof is a computation in interval arithmetic, described in
\Cref{app:cert114}; it is the only step of the paper that depends on a
computer. The point $p^{*}$ is the unique zero of $E(\,\cdot\,,\kappa(p^{*}))$
in an explicit box of radius $10^{-95}$, and the residue functional is
enclosed in a disc of radius $10^{-10}$ around
$75.2385455180-19.2657017116\,i$.

\begin{corollary}\label{cor:family2}
Let $\mathcal P_{1}$ be the irreducible component of $\mathcal P$ through
$p^{*}$, and let $T_{2}\subseteq\mathcal P_{1}$ be the set of points at which
the conditions \textup{(i)} and \textup{(ii)} of \Cref{prop:family2} hold.
Then $T_{2}$ is a smooth connected curve, the forms above are a block family
of character $\alpha_{1}$ over $T_{2}$, $(0,3)$ is a chart, and
$I_{03}(p^{*},v^{*})\ne0$.
\end{corollary}

\begin{proof}
At a point of $\mathcal P$ where $J$ is invertible, $\mathcal P$ is smooth of
dimension one, with $\kappa$ as a local coordinate, by the Jacobian
criterion; in particular only one component passes through it. So $T_{2}$ is
a nonempty Zariski-open subset of the irreducible curve $\mathcal P_{1}$, and
it is smooth and connected. The blocks are $F$, $H$, $G$, $K$ and $t$, with
exponents $(3,1,2,0,0)$, $(0,1,5,0,1)$, $(0,0,0,2,0)$, $(1,5,0,0,1)$ and
$(0,0,0,0,19)$ in $(u_{0},\dots,u_{4})$ and weights
\begin{gather*}
  3+25+2\cdot43=114,\qquad 25+5\cdot43+102=342,\qquad 2\cdot57=114,\\
  1+5\cdot25+102=228,\qquad 19\cdot102=1938=17\cdot114 .
\end{gather*}
Every $u_{k}$ has degree $24$. The forms $F,H,G,K$ are monic in $s$, so none
of them vanishes at $(1:0)$, and with (ii) the blocks are pairwise coprime;
the coefficients $1,a,b,-(1+a+b),c$ do not vanish. This proves (B1)--(B3).
The forms $u_{0}=F^{3}K$ and $u_{3}=-(1+a+b)G^{2}$ share no block, so
$(0,3)$ is a chart, and (iii) is the last assertion.
\end{proof}

The residue functional is computed as follows. Here $A=FH^{3}GK^{2}t^{17}$,
and in the coordinate $s$,
\[
  \eta_{v}=\frac{A\,\Delta_{v}}{u_{0}u_{3}}\,ds=\frac{\mathcal D_{v}}
  {\mathcal M\,F^{5}}\,ds,\qquad \mathcal D_{v}=\det\begin{pmatrix}u_{1}&u_{2}&
  u_{4}\\ \dot u_{1}&\dot u_{2}&\dot u_{4}\\ u'_{1}&u'_{2}&u'_{4}\end{pmatrix},
  \qquad\mathcal M=-abc(1+a+b)\,K^{5}H^{4}G,
\]
since $u_{0}u_{1}u_{2}u_{3}u_{4}=\mathcal M F^{5}\cdot FH^{3}GK^{2}$. The
zeros of $u_{0}$ are the roots of $F$ and of $K$. At a root of $K$ the form
$\eta_{v}$ is regular, by \Cref{lem:etapoles}(a): the block $K$ has weight
$2m$ and occurs once in $u_{0}$ and not in $u_{3}$. As $\mathcal M$ and $F$
are coprime, there is a unique polynomial $R_{v}$ of degree less than $35$
with $R_{v}\mathcal M\equiv\mathcal D_{v}$ modulo $F^{5}$; then
$\eta_{v}-R_{v}F^{-5}ds$ is regular at the roots of $F$, and $R_{v}F^{-5}$
has no other poles in $\CC$ and behaves like $r_{34}s^{-1}$ at infinity,
where $r_{34}$ is the coefficient of $s^{34}$ in $R_{v}$. By the residue
theorem,
\[
  I_{03}(p,v)=r_{34}.
\]
This is the quantity enclosed by the certificate, for $p=p^{*}$ and $v=v^{*}$.

\subsection{Proof of \texorpdfstring{\Cref{thm:fermat114}}{Theorem A}}
\label{ssec:proof114}

\begin{proof}[Proof of \Cref{thm:fermat114}]
\emph{Degree $114$.} By \Cref{prop:family1} at $\lambda=2$ and
\Cref{thm:coniveau}, the character $13\alpha_{2}$ has the property of
\Cref{thm:coniveau}; so has $\alpha_{2}$, as
$M_{13\alpha_{2}}=M_{\alpha_{2}}$. By \Cref{cor:family2} and
\Cref{thm:coniveau}, so has $\alpha_{1}$. By \Cref{lem:twochars}(a) and
\Cref{prop:join}, the multiset $\beta$ of \Cref{lem:twochars}(b) is
algebraic. By \Cref{cor:B114}, $B_{114}=S_{114}+\ZZ[\beta]$, and
\Cref{lem:transfer} gives the Hodge conjecture for $X^{n}_{114}$ for every
$n$, and the algebraicity of every balanced multiset of an even number of
nonzero elements of $\ZZ/114$.

\emph{Degree $57$.} Let $A'$ be a balanced multiset of an even number of
nonzero elements of $\ZZ/57$. Then $2A'\subseteq\ZZ/114$ is balanced, hence
algebraic, and so $A'$ is algebraic by \Cref{lem:algtools}(d). As at the
end of the proof of \Cref{lem:transfer}, this gives the Hodge conjecture for
$X^{n}_{57}$ for every $n$.

\emph{The general case.} Let $m=2^{a}3^{b}19^{c}$ with $a\le1$. For $m\le2$
the variety $X^{n}_{m}$ is a linear space or a quadric, whose cohomology is
spanned by classes of linear subspaces; let $m\ge3$. The elements of
$\mathsf P$ that divide $m$ are among $3$ and $19$, so $\mathcal Q(m)$ is
$\{57\}$ if $57\mid m$ and empty otherwise. By \Cref{thm:aoki},
$B_{m}=S_{m}+\sum_{t}\ZZ\,t\cdot(m/57)_{*}\xi_{57}$, the sum being empty when
$57\nmid m$. The element $\xi_{57}\in B_{57}$ is an integral combination of
classes of balanced multisets with an even number of elements, which are
algebraic by the case of degree $57$; their images under $(m/57)_{*}$ and
then under a unit are classes of algebraic balanced multisets, by
\Cref{lem:algtools}(c),(d). \Cref{lem:transfer} gives the Hodge conjecture
for $X^{n}_{m}$ for every $n$.

\emph{Products and abelian varieties.} By \cite[Theorem~IV]{Shi79}, the
Hodge conjecture for $X^{r_{1}}_{m}\times\dots\times X^{r_{k}}_{m}$ follows
from the Hodge conjecture for $X^{r}_{m}$ for the even $r\le r_{1}+\dots+
r_{k}+2(k-1)$. By \cite[Corollary~3.2]{Aok00}, the Hodge conjecture for
$X^{n}_{m}$ for every $n$ implies it for every abelian variety of Fermat
type of degree $m$.
\end{proof}

\begin{remark}\label{rem:new114}
(a) When $57\nmid m$, \Cref{thm:fermat114} is due to Aoki
\cite[Theorem~0.1]{Aok00}, since $m$ is then of the form $2^{a}3^{b}$,
$19^{c}$ or $2\cdot19^{c}$. For $m=57$, the Hodge conjecture for $X^{n}_{57}$
is stated in \cite[Theorem~1.1]{MMRV26}. Its proof writes $\xi_{57}$, modulo
$S_{114}$, as a balanced sextuple of residues modulo $114$
\cite[Table~4]{MMRV26}, a Hodge class on $X^{4}_{114}$, and declares it
algebraic by Kang's statement \cite[Corollary~3.2]{Kan16}, which rests on the
step discussed in \Cref{rem:kang}; the proof above does not use it. By
Aoki's reduction, the case $m=57$ implies the others, so the same remark
applies to every $m$ divisible by $57$; for $m\ne57$, the first of which are
$114$, $171$, $342$ and $513$, the theorem has not been stated before.
Jumagulov's census of the Fermat fourfolds of even degree at most $250$
lists the degree $114$ as open \cite{Jum26b}.

(b) The census \cite{Jum26b}, a computer-assisted preprint whose lattice
computation we have repeated, finds that for even $m\le250$ every Hodge
class on $X^{4}_{m}$ is algebraic by linear subspaces, Aoki's cycles and the
inductive structure, except for classes attached to the levels $70$, $110$
and $114$ and their multiples $140$, $210$, $220$ and $228$. Peterson's
theorem covers $70$, $140$ and $210$ \cite{Pet26}, and
\Cref{thm:fermat114}, with \Cref{lem:algtools}(d), covers $114$ and $228$.
So the even degrees up to $250$ for which the Hodge conjecture for
$X^{4}_{m}$ remains open are $110$ and $220$, where the classes left over
are those the census attaches to the level $110$. They lie in the class of
$2_{*}\xi_{55}$, and we have not found block families for them.

(c) The characters $\alpha_{1}$, $\alpha_{2}$ and the two families were
found by computer searches over the level-one characters of
$X^{3}_{114}$ and over the possible block structures; the proofs do not
depend on these searches. The searches found no family of degree less than
$24$ for the orbit of $\alpha_{1}$, and no pair of characters as in
\Cref{lem:twochars} for which both families have degree $3$, like the first
one. With such a pair, no step of the proof would need a computer.
\end{remark}
